%%%% RESUMO
%%
%% Apresentação concisa dos pontos relevantes de um texto, fornecendo uma visão rápida e clara do conteúdo e das conclusões do
%% trabalho.

\begin{resumoutfpr}%% Ambiente resumoutfpr
Este trabalho investiga a viabilidade de distribuir uma aplicação prova de conceito baseada no Paradigma Orientado a Notificações (NOP) em um cluster Kubernetes, considerando aspectos como desacoplamento entre componentes, escalabilidade, consumo de recursos e organização da execução em ambiente distribuído. A pesquisa parte do contexto de sistemas com demanda computacional elevada e da necessidade de avaliar se os princípios do NOP, especialmente a redução de redundâncias estruturais e temporais, podem favorecer a construção de uma arquitetura mais adequada à distribuição. Para isso, a fundamentação teórica foi baseada no NOP Query, método orientado a notificações aplicado ao processamento de consultas em fluxo contínuo, cuja proposta demonstra que entidades pequenas e colaborativas podem operar com baixa latência e custo computacional tratável em cenários de grande volume de dados. A partir dessa base, a investigação analisa como os componentes de uma aplicação prova de conceito podem ser conteinerizados e orquestrados pelo Kubernetes, de modo a explorar implantação independente, comunicação entre serviços e possibilidade de expansão horizontal. Os resultados esperados indicam que a combinação entre NOP e Kubernetes é tecnicamente plausível para esse tipo de aplicação, embora ainda existam limites relacionados ao gerenciamento de estado, à comunicação entre componentes e ao consumo de memória em estruturas notificantes. Conclui-se que a distribuição em Kubernetes é compatível com a proposta do NOP e pode ampliar sua aplicabilidade em sistemas modernos, desde que a arquitetura seja planejada com atenção aos custos de coordenação e às restrições da prova de conceito.

\end{resumoutfpr}

